\documentclass[11pt,a4paper,sans]{moderncv}

\moderncvstyle{banking}
\moderncvcolor{blue}

\usepackage[utf8]{inputenc}
\usepackage{lmodern}
\usepackage[scale=0.85]{geometry}
\renewcommand*{\namefont}{\fontsize{28}{32}\mdseries\upshape}
\renewcommand*{\titlefont}{\fontsize{22}{26}\mdseries\upshape}

% Personal Information
\name{Ibrahim}{Sayar}
\title{Senior Software Engineer}
\address{Augsburg, Deutschland}{}{}
\phone[mobile]{+49~170~1030979}
\email{ibrahimsayar96@gmail.com}
\homepage{ibrahimsayar.com}
\social[linkedin]{ibrahimsayar}
\social[github]{ibrahimsayar}
\extrainfo{Arbeitserlaubnis: Deutschland (kein Sponsoring erforderlich)}

\begin{document}

\makecvtitle

\vspace{-2em}

%-----------------------------------------------------------------------
% ZUSAMMENFASSUNG
%-----------------------------------------------------------------------
\section{\"Uber mich}
Senior Software Engineer mit \"uber 10 Jahren Erfahrung in Softwarearchitektur, Backend-Entwicklung und DevOps. Ich entwerfe und betreibe produktive Systeme von Ende zu Ende: Service- und Integrationsarchitektur, PHP/Laravel- und Node.js/NestJS-Services, CI/CD, Cloud-Infrastruktur und Monitoring. Aktuell verantworte ich die Architektur der E-Commerce-Plattformen mehrerer Marken und unterst\"utze Spiele und weitere Produkte. Besonders gern modernisiere ich Legacy-Systeme Schritt f\"ur Schritt, ohne den laufenden Betrieb zu unterbrechen.

%-----------------------------------------------------------------------
% BERUFSERFAHRUNG
%-----------------------------------------------------------------------
\section{Berufserfahrung}

\cventry{Mai 2026 -- Heute}{Senior Software Engineer}{Naylalabs}{Remote}{}{
\begin{itemize}
\item Verantwortung f\"ur Softwarearchitektur und DevOps der E-Commerce-Plattformen mehrerer unabh\"angiger Marken mit zusammen 300.000+ Bestellungen pro Monat \"uber Web, Mobile Apps und Marktpl\"atze; Architektur- und Infrastruktur-Support f\"ur Spiele und weitere Produkte; Team von 25+ Personen.
\item Festlegung von Architektur und Technologie je Produkt (PHP/Laravel, NestJS, React); Entwurf ereignisgesteuerter, queue-basierter Integrationen (DDD, BullMQ/Redis) mit Marktpl\"atzen und ERP-Systemen sowie schrittweise Abl\"osung von Legacy-Integrationen ohne Ausfallzeit.
\item Betrieb der Produktionsinfrastruktur: Linux, nginx, Load Balancing, Cloudflare (500 Mio.+ Requests pro Monat), PostgreSQL/MySQL und Monitoring mit Grafana; Leitung von Datenbankmigrationen.
\end{itemize}
}

\cventry{Jan 2025 -- Mai 2026}{TestOps Engineer}{Insider One}{T\"urkei}{}{
\begin{itemize}
\item Aufbau und Pflege automatisierter Testpipelines (Selenium, Appium, Postman) mit CI-Integration auf Kubernetes f\"ur eine Plattform mit 7 Mrd.+ Events pro Woche.
\item Zusammenarbeit mit Entwicklungsteams zur Verbesserung von Release-Qualit\"at und Zuverl\"assigkeit.
\end{itemize}
}

\cventry{Aug 2023 -- Feb 2025}{DevOps Engineer}{PangeaVR}{Istanbul, T\"urkei}{}{
Verwaltung der Cloud-Infrastruktur auf AWS mit Docker- und Kubernetes-Orchestrierung.
Aufbau und Pflege von CI/CD-Pipelines f\"ur automatisierte Bereitstellungen.
Administration von Linux-Servern (CentOS) und Optimierung der Systemleistung.
}

\cventry{Aug 2022 -- Aug 2023}{Full-stack-Entwickler \& DevOps}{Naylalabs}{Samsun, T\"urkei}{}{
Entwicklung von Full-Stack-Webanwendungen mit Node.js, Python und JavaScript.
Entwurf von RESTful-APIs und Verwaltung von PostgreSQL/MySQL-Datenbanken.
Implementierung containerisierter Bereitstellungen mit Docker und CI/CD-Automatisierung.
}

\cventry{Jul 2021 -- Sep 2022}{Webentwickler}{Naylalabs}{Samsun, T\"urkei}{}{
Entwicklung responsiver Webanwendungen mit modernen JavaScript-Frameworks.
Integration von Drittanbieterdiensten einschlie\ss{}lich E-Mail-Services und JSON-APIs.
Zusammenarbeit in agilen Arbeitsabl\"aufen mit Git-Versionskontrolle.
}

\cventry{Jan 2020 -- Jun 2021}{Leitender Back-end-Entwickler}{Veriweb}{Samsun, T\"urkei}{}{
Leitung des Back-end-Entwicklungsteams, Definition von Softwarearchitektur und Programmierstandards.
Entwicklung skalierbarer serverseitiger Anwendungen mit PHP und Python.
Entwurf und Optimierung von SQL-Datenbanken (MySQL, PostgreSQL).
}

\cventry{Nov 2016 -- Dez 2019}{Mitgr\"under \& Back-end-Entwickler}{Kodakil}{Samsun, T\"urkei}{}{
Mitgr\"undung eines Tech-Startups, Leitung der Back-end-Architektur und Entwicklung.
Aufbau von RESTful-APIs und Verwaltung der Serverinfrastruktur unter Linux.
\"Uberwachung der technischen Strategie und des Produktentwicklungszyklus.
}

\cventry{Mai 2014 -- Nov 2016}{Webdesign, Programmierung \& IT}{Fikir Kulubu}{Samsun, T\"urkei}{}{
Entwicklung und Wartung von Webanwendungen mit PHP und JavaScript.
Verwaltung der IT-Infrastruktur und technischer Support.
}

\cventry{Feb 2013 -- Feb 2014}{Webentwickler}{Teknocenter}{}{}{
Entwicklung von Websites mit HTML, CSS, JavaScript und PHP.
}

%-----------------------------------------------------------------------
% TECHNISCHE KENNTNISSE
%-----------------------------------------------------------------------
\section{Technische Kenntnisse}

\cvitem{Architektur}{Domain-Driven Design, Microservices, ereignisgesteuerte \& queue-basierte Systeme, Systemintegration (Marktpl\"atze, ERP), Legacy-Modernisierung, Caching, Load Balancing}
\cvitem{Sprachen}{PHP, TypeScript, JavaScript, Python, SQL, Bash}
\cvitem{Back-end}{Laravel, Node.js, NestJS, REST-APIs, BullMQ, Redis}
\cvitem{Front-end}{React, Vite, HTML, CSS}
\cvitem{DevOps \& Cloud}{GitHub Actions, Docker, Kubernetes, Terraform, AWS, Linux, nginx, Cloudflare, CI/CD}
\cvitem{Monitoring \& Sicherheit}{Grafana, OAuth 2.0/JWT, Cloudflare WAF, SSL/TLS}
\cvitem{Datenbanken}{PostgreSQL, MySQL, Oracle (Integration)}
\cvitem{Tests}{Selenium, Appium, Postman, E2E-Automatisierung}

%-----------------------------------------------------------------------
% AUSBILDUNG
%-----------------------------------------------------------------------
\section{Ausbildung}

\cventry{2017 -- 2023}{Bachelor -- Wirtschaftsinformatik}{Anadolu Universit\"at}{}{}{Zeugnisbewertung (ZAB): beantragt}
\cventry{2014 -- 2016}{Associate Degree -- Informatik}{Ondokuz Mayis Universit\"at}{}{}{}
\cventry{2011 -- 2014}{Informatik}{Atakum Technische und Industrielle Berufsschule}{}{}{}

%-----------------------------------------------------------------------
% SPRACHEN
%-----------------------------------------------------------------------
\section{Sprachen}

\cvitemwithcomment{T\"urkisch}{Muttersprache}{}
\cvitemwithcomment{Englisch}{B1 (Mittelstufe)}{}
\cvitemwithcomment{Deutsch}{A2 (Integrationskurs, laufend)}{}

%-----------------------------------------------------------------------
% AUSZEICHNUNGEN
%-----------------------------------------------------------------------
\section{Auszeichnungen}

\cvitem{}{Ondokuz Mayis Universit\"at -- Urkunde mit Auszeichnung}

\end{document}
